% Generated by tools/edn_latex.py from reports/edn.md.
% Do not edit: change reports/edn.md and run `make edn`.
\ednentry{
  id = {EDN-46},
  title = {DagsHub credentials live in two gitignored stores, and the token is entered twice},
  date = {2026-09-30},
  milestone = {M2: Reproducibility},
  topic = {Credential handling, developer onboarding},
  participants = {Lukas},
  decision = {A contributor's DagsHub token is stored twice, in two gitignored files: in \texttt{.\allowbreak{}env} as \texttt{MLFLOW\_\allowbreak{}TRACKING\_\allowbreak{}PASSWORD}, which MLflow reads from the environment, and in \texttt{.\allowbreak{}dvc/\allowbreak{}config.\allowbreak{}local} via \texttt{dvc remote modify origin -\allowbreak{}-\allowbreak{}local password}, which is the only place DVC's HTTP remote reads a password from.
Nothing derives one store from the other, and \texttt{.\allowbreak{}env.\allowbreak{}template} deliberately does not list \texttt{DAGSHUB\_\allowbreak{}USERNAME} or \texttt{DAGSHUB\_\allowbreak{}USER\_\allowbreak{}TOKEN}.
\texttt{recommenditos/\allowbreak{}config.\allowbreak{}py} loads \texttt{<repo>/\allowbreak{}.\allowbreak{}env} by name rather than searching upwards, so the credentials of a checkout are its own.},
  rationale = {\emph{Alternatives considered:}
\begin{itemize}[leftmargin=*, topsep=2pt]
\item \textbf{Option A (chosen): two stores, the token pasted into each, and a template that lists only what the project reads.}
\newline Pros: no machinery to maintain; both stores are already gitignored, so no token can reach a commit; each tool reads its credentials the way it documents; the template cannot mislead, because every variable in it is read and a test asserts that.
\newline Cons: the token is entered twice, and the reader has to be told why or it looks like an oversight; the \texttt{dvc remote modify} form puts the token in shell history and in \texttt{ps}.
\item \textbf{Option B: one source of truth in \texttt{.\allowbreak{}env}, plus a script that writes \texttt{.\allowbreak{}dvc/\allowbreak{}config.\allowbreak{}local} from it.}
\newline Pros: one paste; one file to rotate.
\newline Cons: a piece of project machinery whose whole job is to write a secret to disk on its own initiative, which is a worse thing to own than a second paste; another step that can fail or drift; nothing in the course requires it.
\item \textbf{Option C: \texttt{.\allowbreak{}env} only, relying on an environment variable for DVC.}
\newline Pros: would be the simplest of all, if it existed.
\newline Cons: it does not. Measured rather than assumed: with \texttt{MLFLOW\_\allowbreak{}*}, \texttt{DAGSHUB\_\allowbreak{}USERNAME} and \texttt{DAGSHUB\_\allowbreak{}USER\_\allowbreak{}TOKEN} all exported, \texttt{dvc status -\allowbreak{}c} still fails with \texttt{configuration error -\allowbreak{} HTTP \textquotesingle{}basic\textquotesingle{} authentication require both \textquotesingle{}user\textquotesingle{} and \textquotesingle{}password\textquotesingle{}}, and DVC 3.67.1's HTTP remote schema offers exactly \texttt{user}, \texttt{password} and \texttt{ask\_\allowbreak{}password} and no environment route at all.
\item \textbf{Option D: \texttt{ask\_\allowbreak{}password true}, prompting instead of storing.}
\newline Pros: nothing on disk; nothing in shell history.
\newline Cons: \texttt{dvc\_\allowbreak{}http} calls \texttt{getpass}, which needs a terminal, so it prompts on every \texttt{dvc pull} and breaks any unattended \texttt{dvc repro} or CI job. Documented as the better choice on a shared machine, not as our default.
\end{itemize}
\par
\emph{Why:} Option C would have been the right answer and is not available, which is what makes the double entry a property of DVC rather than a choice.
Between A and B, the property that matters is that no token can reach a commit, and both stores already have it; B buys one fewer paste at the price of owning a secret-writing script.
Listing \texttt{DAGSHUB\_\allowbreak{}*} in the template was actively harmful: nothing reads those names, so filling them in configures nothing while reading as though it configured DVC, which is the confusion the getting-started page exists to prevent.
Scoping \texttt{.\allowbreak{}env} to the repository belongs to the same decision: a bare \texttt{load\_\allowbreak{}dotenv()} searches upwards, so a clone nested under another checkout inherited that one's token and looked configured when it was not.},
  ai = {Information seeking, Alternative assessment, Solution generation},
  response = {Accepted with modifications},
  assessment = {AI established by execution that DVC has no environment-variable route, which is the fact the whole decision rests on, and it wrote the two-store walkthrough. An adversarial review of that work, also by AI, then found three things the first pass had asserted rather than checked: that the \texttt{.\allowbreak{}env} scoping was not repository-local, so the review's own ``fresh clone'' verification had in fact been running on the parent checkout's credentials; that \texttt{mlflow.\allowbreak{}db} was neither gitignored nor prevented; and that the documented newcomer command printed a traceback where the page promised a message. All three were reproduced before being fixed. The modification is that the template lost the \texttt{DAGSHUB\_\allowbreak{}*} block, which the first pass had defended as documentation, once the review showed nothing reads it.},
  aievidence = {Claude Code sessions on 2026-09-30: the DVC environment-variable question was settled by running \texttt{dvc status -\allowbreak{}c} with and without the variables exported; the review findings were each reproduced before any fix, including a probe showing \texttt{find\_\allowbreak{}dotenv} resolving to a \texttt{.\allowbreak{}env} three directories above a credential-free worktree.},
  evidence = {\href{https://github.com/mlops-2627q1-mds-upc/MLOPS_Recommenditos/issues/42}{issue \#42}; \href{https://github.com/mlops-2627q1-mds-upc/MLOPS_Recommenditos/pull/50}{PR \#50}; \hyperref[EDN-19]{EDN-19}; NFR-15 in \href{https://github.com/mlops-2627q1-mds-upc/MLOPS_Recommenditos/blob/main/docs/docs/specification.md}{the specification}, which held this as part of NFR-09 when this entry was written (\hyperref[EDN-66]{EDN-66}).},
}
